\section{Hydrogen atom as a minimal proton--electron composite (interface + matching + audit)}
\label{app:hydrogen_atom_interface}

\paragraph{Scope and layer discipline.}
\InterfaceTag This appendix records a minimal interface construction of the hydrogen atom as a composite system built from the hydrogen nucleus (proton) and an electron.
The goal is to make the composite structure and the matching targets explicit while respecting the paper's scope boundaries: we do not claim a theorem-level derivation of atomic spectra from the $m=6$ folding core.

\subsection{Composite-system structure (interface)}
\paragraph{Tensor-product packaging.}
\InterfaceTag We treat the hydrogen atom as a two-subsystem composite with Hilbert space
\[
\mathcal{H}_{\mathrm{H}} \;:=\; \mathcal{H}_p\otimes \mathcal{H}_e,
\]
and joint readouts organized by product POVMs.
This uses the minimal composite-system package recorded in Appendix~\ref{app:composite_systems_tensor_products}.

\paragraph{Carrier dictionary (interface pointer).}
\InterfaceTag The proton is treated as a minimal color-singlet composite carrier (Appendix~\ref{app:proton_hydrogen_nucleus_interface}), and the electron is a minimal matter anchor at $m=6$ under the closed SM labeling dictionary (Part~IV).
The hydrogen atom is then a composite carrier whose low-energy structure (binding) is recorded at the matching layer.

\subsection{Vendored reference targets (matching/audit)}
\label{subsec:hydrogen_atom_miniset}
\MatchTag Table~\ref{tab:hydrogen_atom_miniset} records a tiny vendored reference set for hydrogen-atom targets (Bohr radius and representative energy scales), used only as matching-layer conventions.

\AuditTag \paragraph{Audit summary (hydrogen atom miniset).} \AuditTag Rows are vendored matching-layer reference targets for the hydrogen atom. The JSON input is stored under \texttt{data/pdg\_minisets/hydrogen\_atom\_miniset.json}; this fragment is generated by \texttt{scripts/exp\_hydrogen\_atom\_miniset.py}.%

\begin{table}[t]
\centering
\small
\setlength{\tabcolsep}{8pt}
\renewcommand{\arraystretch}{1.15}
\begin{tabular}{llll}
\toprule
quantity & value & unit & source \\
\midrule
bohr\_radius & 5.291772109e-11 & m & CODATA (Bohr radius, representative) \\
ionization\_energy\_n1 & 13.598434 & eV & NIST/CODATA/PDG-style reference (representative) \\
rydberg\_energy & 13.60569312 & eV & CODATA (Rydberg energy, representative) \\%
\bottomrule
\end{tabular}
\caption{Vendored hydrogen-atom reference targets (matching layer). Rows are generated by \texttt{scripts/exp\_hydrogen\_atom\_miniset.py} from \texttt{data/pdg\_minisets/hydrogen\_atom\_miniset.json}.}
\label{tab:hydrogen_atom_miniset}
\end{table}

\subsection{Resolution mapping for the binding-energy scale (audit)}
\label{subsec:hydrogen_binding_resolution_mapping}
\AuditTag The hydrogen ground-state ionization energy is a \emph{composite binding scale}.
Accordingly, mapping it to the staircase label $m_{\mathrm{eff}}$ is reported only as a scale label in the $r(\mu)$ coordinate, not as a claim that single-window localized matter exists for $m<6$ (cf.\ the localization admissibility threshold in Appendix~\ref{app:inference_ledger}).

\AuditTag \paragraph{Audit summary (hydrogen binding scale $\mu\to r\to m_{\mathrm{eff}}$).} \AuditTag We map the ground-state ionization energy (a composite binding scale) from the vendored miniset to $\mu$ in GeV and then to the resolution coordinate $r(\mu)=\log(\mu/m_e)/\log\varphi$ using $m_e=\texttt{M\_E\_GEV}$ and $\varphi=\texttt{PHI}$. We report the staircase label $m_{\mathrm{eff}}(\mu)=6+\lfloor r(\mu)/r_{\mathrm{step}}\rfloor$ with $r_{\mathrm{step}}=2\pi$ (Section~\ref{sec:falsifiability_predictions}). Interpretation boundary: values with $m_{\mathrm{eff}}<6$ are \emph{not} treated as single-window localized matter (Section~\ref{subsec:matter_as_types}); they are recorded only as scale labels for composite low-energy structure.%

\begin{table}[t]
\centering
\small
\setlength{\tabcolsep}{6pt}
\renewcommand{\arraystretch}{1.15}
\begin{tabular}{lrrrrrrl}
\toprule
object & $E$ [eV] & $\mu$ [GeV] & $r(\mu)$ & $m_{\mathrm{eff}}$ & $\mu_{\mathrm{th}}(m_{\mathrm{eff}})$ [GeV] & $\mu_{\mathrm{th}}(m_{\mathrm{eff}}+1)$ [GeV] & note \\
\midrule
hydrogen (n=1 ionization) & 13.598434 & 1.35984\times 10^{-8} & -21.8909 & 2 & 2.85751\times 10^{-9} & 5.87619\times 10^{-8} & composite binding scale (below m=6 single-window admissibility) \\%
\bottomrule
\end{tabular}
\caption{Resolution mapping for a hydrogen composite binding-energy scale under the fixed staircase convention $r_{\mathrm{step}}=2\pi$ (Section~\ref{sec:falsifiability_predictions}). Rows are generated by \texttt{scripts/exp\_hydrogen\_binding\_resolution\_mapping.py}.}
\label{tab:hydrogen_binding_resolution_mapping}
\end{table}

\subsection{Hydrogen spectral-line targets (matching/audit)}
\label{subsec:hydrogen_spectral_lines_miniset}
\MatchTag Table~\ref{tab:hydrogen_spectral_lines_miniset} records a tiny vendored reference set for hydrogen spectral-line targets (including the 21cm hyperfine line), used only as matching-layer conventions.

\AuditTag \paragraph{Audit summary (hydrogen spectral-line miniset).} \AuditTag Rows are vendored matching-layer reference targets for hydrogen spectral lines (e.g.\ 21cm hyperfine). The JSON input is stored under \texttt{data/pdg\_minisets/hydrogen\_spectral\_lines\_miniset.json}; this fragment is generated by \texttt{scripts/exp\_hydrogen\_spectral\_lines\_miniset.py}.%

\begin{table}[t]
\centering
\small
\setlength{\tabcolsep}{7pt}
\renewcommand{\arraystretch}{1.15}
\begin{tabular}{llll}
\toprule
quantity & value & unit & source \\
\midrule
hyperfine\_21cm\_frequency & 1420405752 & Hz & NIST/IAU-style reference (representative): hydrogen hyperfine transition frequency \\
hyperfine\_21cm\_wavelength & 0.2110611405 & m & Derived reference (representative): lambda=c/f using the same frequency convention \\
hyperfine\_21cm\_energy & 5.87433e-06 & eV & Derived reference (representative): E=h f in eV \\
lyman\_alpha\_wavelength & 1.21567e-07 & m & NIST-style reference (representative): Lyman-alpha vacuum wavelength \\
lyman\_alpha\_frequency & 2.46607e+15 & Hz & Derived reference (representative): f=c/lambda using exact c and the vendored wavelength \\
lyman\_alpha\_energy & 10.1988 & eV & Derived reference (representative): E=hc/lambda in eV \\%
\bottomrule
\end{tabular}
\caption{Vendored hydrogen spectral-line reference targets (matching layer). Rows are generated by \texttt{scripts/exp\_hydrogen\_spectral\_lines\_miniset.py} from \texttt{data/pdg\_minisets/hydrogen\_spectral\_lines\_miniset.json}.}
\label{tab:hydrogen_spectral_lines_miniset}
\end{table}

\subsection{Frequency-first bridge to the phase$\to$delay audits (interface/audit)}
\label{subsec:hydrogen_phase_delay_bridge}
\InterfaceTag Spectral lines provide characteristic frequencies. In the frequency-first interface, a line at frequency $f$ defines $\omega=2\pi f$.
Whenever a platform supplies a complex response $S(\omega)$ (or transfer function) whose phase $\phi(\omega)=\arg S(\omega)$ is measurable on a band, the manuscript's standard phase$\to$delay step applies:
in one channel, $\tau(\omega)=\dd\phi/\dd\omega$ (Appendix~\ref{app:force_phase_delay_audit}), and the same numerical-stability discipline can be applied to discrete phase samples (Appendix~\ref{app:k4_scattering_phase_delay_audit}).
\AuditTag A dedicated spectral phase$\to$delay audit (with the same window-family stability sweep form) is recorded in Appendix~\ref{app:k4_spectral_phase_delay_audit}.
\AuditTag This is an interface bridge statement only; it does not assert that atomic lines are scattering phase shifts.

\AuditTag \paragraph{Audit summary (hydrogen spectral frequency$\to\omega$ bridge).} \AuditTag For a spectral line with frequency $f$, define $\omega=2\pi f$ and period $T=1/f$. Once a platform supplies a response phase $\phi(\omega)$ (e.g.\ via a complex transfer function), the paper's standard phase$\to$delay step applies: $\tau(\omega)=\dd\phi/\dd\omega$ (see Appendix~\ref{app:force_phase_delay_audit} and Appendix~\ref{app:k4_scattering_phase_delay_audit}). Rows are generated by \texttt{scripts/exp\_hydrogen\_spectral\_phase\_delay\_bridge.py} from \texttt{data/pdg\_minisets/hydrogen\_spectral\_lines\_miniset.json}.%

\begin{table}[t]
\centering
\small
\setlength{\tabcolsep}{6pt}
\renewcommand{\arraystretch}{1.15}
\begin{tabular}{lrrrr}
\toprule
quantity & $f$ [Hz] & $\omega=2\pi f$ [rad/s] & $T=1/f$ [s] & note \\
\midrule
hyperfine\_21cm\_frequency & 1.42041\times 10^{9} & 8.92467\times 10^{9} & 7.04024\times 10^{-10} & hyperfine (phase/delay bridge via omega) \\
lyman\_alpha\_frequency & 2.46607\times 10^{15} & 1.54948\times 10^{16} & 4.05503\times 10^{-16} & UV transition (phase/delay bridge via omega) \\%
\bottomrule
\end{tabular}
\caption{Frequency-first bridge table for hydrogen spectral lines. Rows are generated by \texttt{scripts/exp\_hydrogen\_spectral\_phase\_delay\_bridge.py}.}
\label{tab:hydrogen_spectral_delay_bridge}
\end{table}

\subsection{Resolution mapping for the 21cm hyperfine scale (audit)}
\label{subsec:hydrogen_21cm_resolution_mapping}
\AuditTag The 21cm hyperfine splitting is a composite low-energy scale (internal structure of the hydrogen ground state).
Its $m_{\mathrm{eff}}$ value is reported only as a scale label in the $r(\mu)$ coordinate, not as a single-window admissibility claim.

\AuditTag \paragraph{Audit summary (hydrogen 21cm $\mu\to r\to m_{\mathrm{eff}}$).} \AuditTag We map the 21cm hyperfine splitting energy (a composite low-energy scale) from the vendored miniset to $\mu$ in GeV and then to $r(\mu)=\log(\mu/m_e)/\log\varphi$ using $m_e=\texttt{M\_E\_GEV}$. We report the staircase label $m_{\mathrm{eff}}(\mu)=6+\lfloor r(\mu)/r_{\mathrm{step}}\rfloor$ with $r_{\mathrm{step}}=2\pi$ (Section~\ref{sec:falsifiability_predictions}). Interpretation boundary: $m_{\mathrm{eff}}<6$ is recorded only as a scale label for a composite effect.%

\begin{table}[t]
\centering
\small
\setlength{\tabcolsep}{6pt}
\renewcommand{\arraystretch}{1.15}
\begin{tabular}{lrrrrrrl}
\toprule
object & $E$ [eV] & $\mu$ [GeV] & $r(\mu)$ & $m_{\mathrm{eff}}$ & $\mu_{\mathrm{th}}(m_{\mathrm{eff}})$ [GeV] & $\mu_{\mathrm{th}}(m_{\mathrm{eff}}+1)$ [GeV] & note \\
\midrule
hydrogen 21cm hyperfine & 5.87433e-06 & 5.87433\times 10^{-15} & -52.345 & -3 & 7.77047\times 10^{-16} & 1.59792\times 10^{-14} & composite hyperfine scale (below m=6 single-window admissibility) \\%
\bottomrule
\end{tabular}
\caption{Resolution mapping for the hydrogen 21cm hyperfine energy scale under the fixed staircase convention $r_{\mathrm{step}}=2\pi$ (Section~\ref{sec:falsifiability_predictions}). Rows are generated by \texttt{scripts/exp\_hydrogen\_21cm\_resolution\_mapping.py}.}
\label{tab:hydrogen_21cm_resolution_mapping}
\end{table}

\subsection{Resolution mapping for the Lyman-\texorpdfstring{$\alpha$}{alpha} scale (audit)}
\label{subsec:hydrogen_lyman_alpha_resolution_mapping}
\AuditTag The Lyman-$\alpha$ transition energy is a composite atomic transition scale.
Its $m_{\mathrm{eff}}$ value is reported only as a scale label in the $r(\mu)$ coordinate, not as a single-window admissibility claim.

\AuditTag \paragraph{Audit summary (hydrogen Lyman-$\alpha$ $\mu\to r\to m_{\mathrm{eff}}$).} \AuditTag We map the vendored Lyman-$\alpha$ transition energy (a composite atomic scale) to $\mu$ in GeV and then to $r(\mu)=\log(\mu/m_e)/\log\varphi$ using $m_e=\texttt{M\_E\_GEV}$. We report $m_{\mathrm{eff}}(\mu)=6+\lfloor r(\mu)/r_{\mathrm{step}}\rfloor$ with $r_{\mathrm{step}}=2\pi$ (Section~\ref{sec:falsifiability_predictions}). Interpretation boundary: $m_{\mathrm{eff}}<6$ is a scale label only.%

\begin{table}[t]
\centering
\small
\setlength{\tabcolsep}{6pt}
\renewcommand{\arraystretch}{1.15}
\begin{tabular}{lrrrrrrl}
\toprule
object & $E$ [eV] & $\mu$ [GeV] & $r(\mu)$ & $m_{\mathrm{eff}}$ & $\mu_{\mathrm{th}}(m_{\mathrm{eff}})$ [GeV] & $\mu_{\mathrm{th}}(m_{\mathrm{eff}}+1)$ [GeV] & note \\
\midrule
hydrogen Lyman-alpha & 10.1988 & 1.01988\times 10^{-8} & -22.4888 & 2 & 2.85751\times 10^{-9} & 5.87619\times 10^{-8} & composite atomic transition scale (below m=6 single-window admissibility) \\%
\bottomrule
\end{tabular}
\caption{Resolution mapping for the hydrogen Lyman-$\alpha$ energy scale under the fixed staircase convention $r_{\mathrm{step}}=2\pi$ (Section~\ref{sec:falsifiability_predictions}). Rows are generated by \texttt{scripts/exp\_hydrogen\_lyman\_alpha\_resolution\_mapping.py}.}
\label{tab:hydrogen_lyman_alpha_resolution_mapping}
\end{table}
